% Phase 11 results, System 2 (round 1). Numbers from harness_ab_gen.tex.
\textbf{System~2.} System~2 is the code agent's default model, and three tasks
were run per arm (a System~2 step takes minutes on this host). The Phase~11 arm
verified all \HABTwoNewVerified{} in \HABTwoNewCalls{} model calls, one per task,
in \HABTwoNewWall\,s in all; the Phase~10 loop verified
\HABTwoOldVerified{} of \HABTwoOldN{} in \HABTwoOldCalls{} calls and
\HABTwoOldWall\,s. Two of its runs ended at the 900\,s cap, which now holds
inside a model call (both ended at 900.0\,s; in Phase~10 one call had run to
1{,}041\,s), and the third spent its six steps. \HABTwoOldMalformed{} of its
\HABTwoOldActions{} steps were malformed: update hunks whose lines lacked the
\texttt{+}/\texttt{-}/space prefix or the \texttt{@@} marker, and replies with no
action at all, two of them cut off by the token limit. Under the grammar no step
was malformed. Two of the three Phase~11 steps are worth a note: the files were
short, so \texttt{"auto"} offered only a whole file, and the model wrote a Codex
patch \emph{inside} the \verb|```python| block, which the grammar allows (a code
line may start with \texttt{***}); the parser applies a patch envelope wherever it
appears, both applied, and both runs verified. Every Phase~11 call spent its full
reasoning budget (511 tokens): decoding took \HABTwoNewDecodeS\,s of the
\HABTwoNewWall\,s and prefill \HABTwoNewPrefillS\,s, which is why round~2
measures the budget. With three tasks this is a demonstration rather than an
estimate, but the direction agrees with System~1 and the mechanism is visible in
every step.

% Phase 11 results, rounds 2 and 3. Numbers from harness_ab_macros.tex.
\textbf{Round 2: breaking repetition, and the reasoning budget.} Round~1 showed
a failure the grammar cannot touch: a small model that rewrites the same failing
file step after step, anchored on its own previous reply (System~1 reproduced
an already-judged file in four consecutive steps of one run). Two changes went
into the build for round~2. Failing asserts now carry the value they got
(\texttt{assert f(x) == y -> AssertionError: got 5}), the concise execution
feedback that self-debugging relies on \cite{chen2024selfdebug}; and
\texttt{unstick}: after a step that reproduced files the verifier had already
judged, the transcript is compacted before the next call (the model no longer
sees, and copies, its repeated reply; the attempts list says what was tried) and
the step is sampled 0.3 hotter per repeat. Table~\ref{tab:harnessab2} gives the
interleaved comparison on the same ten System~1 tasks: with \texttt{unstick} the
agent verified \HABROneUVerified{} in \HABROneUCalls{} calls and \HABROneUWall\,s,
without it \HABROneNoUVerified{} in \HABROneNoUCalls{} calls and
\HABROneNoUWall\,s; the two discordant tasks both favour \texttt{unstick}, so it is
on by default. The value-carrying evidence was in both arms and is not isolated;
the Phase~11 arm verified 4 of 10 in round~1 and \HABROneNoUVerified{} without
\texttt{unstick} in round~2, which is consistent with it helping but is not an
interleaved comparison.

On System~2 every call of round~1 spent its full reasoning budget, so round~2
compared 512 and 256 reasoning tokens per step on six tasks. Both verified all
six. Halving the budget halved the time of the four tasks it solved in one call
(about 55\,s against 100\,s), but two tasks then needed four and six calls, \HABRTwoSmallMalformed{}
of its \HABRTwoSmallActions{} steps were malformed, and the arm took
\HABRTwoSmallWall\,s against \HABRTwoBigWall\,s: a smaller budget is not faster
overall. The malformed steps exposed a gap in the grammar. With a short file only
a whole file is offered, and System~2 wrote a Codex patch \emph{inside} the
\verb|```python| block, which the grammar allowed (round~1's two such patches
applied; round~2's four did not). A code line may no longer start with
\texttt{*} (a top-level Python line never does), so the edit policy is now
enforced by the grammar, not only by the prompt.

% Phase 11 results, round 3. Numbers from harness_ab_macros.tex.
\textbf{Round 3: a cheap first attempt.} With patches kept out of code blocks,
round~3 compared the 512-token budget on every step with 256 tokens on the
first step and 512 after a failure, on the same six System~2 tasks
(Table~\ref{tab:harnessab3}). Both arms verified all six, each in one call per
task, so the escalation was never needed; the cheap first attempt took
\HABRRTwoAdaptWall\,s in all against \HABRRTwoFixedWall\,s, half the time, and
decoded \HABRRTwoAdaptDecoded{} tokens against \HABRRTwoFixedDecoded. It is the
default (\texttt{first\_step\_thinking\_budget = 256}). The same build primes
llama-server's cache with the agent's system prompt at startup: the first call
of round~3 found it cached (207 tokens) where round~1's first call found 5, and
over the round \HABRRTwoAdaptHit\,\% of the prompt tokens of the faster arm came
from the cache (partly because both arms send the same task message, so the
second arm of a task finds it cached); prefill was \HABRRTwoAdaptPrefillS\,s of its
\HABRRTwoAdaptWall\,s. On this machine a verified System~2 fix of a small
function now takes 33--59\,s, against single calls of 500--1{,}041\,s and no
verified task in the Phase~10 live run; the hardware and the models are the same.


